\chapter{Analiza problemu i~wymagania}
\label{chap:prob_and_req}
% TODO: change after the end of this chapter
Optymalizacja architektury sieci jest złożonym zagadnieniem i~wymaga dokładnego rozplanowania co jest jej poddawane.


Biblioteka udostępnia: funkcję umożliwiającą automatyczne wyszukiwanie optymalnej architektury sieci MLP z wykorzystaniem algorytmu genetycznego oraz funkcję tworzącą sieć z osobnika. Pierwsza funkcja odpowiada za przeprowadzenie całego procesu optymalizacji - od wygenerowania populacji początkowej, poprzez ewolucję kolejnych generacji, aż do zwrócenia wyników i zapisania przebiegu eksperymentu. Przyjmuje ona następujące parametry obligatoryjne: dane treningowe i testowe oraz ścieżkę do zapisu przebiegu optymalizacji. Oraz ma jeden parametr opcjonalny: dane walidacyjne.
Dodatkowo można zmieniać następujące parametry domyślne: liczba atrybutów wejściowych, wielkość populacji, liczba generacji, współczynnik krzyżowania i mutacji, część korygująca funkcji przystosowania, maksymalną liczbę zrównolegleń, liczba obliczań przystosowania, flaga wyświetlania postępu oraz flaga rysowania wykresu. Funkcja zwraca wartości funkcji przystosowania w kolejnych generacjach wraz z populacją końcową.
Przyjmowane dane muszą być w następujących formatach:
\begin{nospitemize}
  \item dane treningowe - tabela danych zawierająca w jednym wierszu atrybuty wchodzące i wychodzące,
  \item dane walidacyjne - tabela danych zawierająca w jednym wierszu atrybuty wchodzące i wychodzące,
  \item dane testowe - tabela danych zawierająca w jednym wierszu atrybuty wchodzące i wychodzące,
  \item ścieżka do zapisu przebiegu optymalizacji - ścieżka do pustego (istniejącego lub nie) folderu,
  \item liczba atrybutów wejściowych - liczba całkowita różna zeru (liczba dodatnia bezpośrednio wskazuje zadaną wartość, a wartość absolutna ujemnej oznacza ile od wszystkich atrybutów należy odjąć by uzyskać zadaną wartość),
  \item wielkość populacji - liczba całkowita dodatnia,
  \item liczba generacji - liczba całkowita dodatnia,
  \item współczynnik krzyżowania - liczba rzeczywista z przdziału $\langle 0, 1 \rangle$,
  \item współczynnik mutacji - liczba rzeczywista z przdziału $\langle 0, 1 \rangle$,
  \item część korygująca funkcji przystosowania - funkcja przyjmująca następujące parametry:
  \begin{nospitemize}
    \item osobnik - osobnik deterministycznie określający strukturę sieci neuronowej,
    \item dopasowanie reakcji nauczonej sieci do przewidywanych odpowiedzi - liczba rzeczywista
  \end{nospitemize}%
  a zwracająca przystosowanie osobnika - liczba rzeczywista wzrastająca wraz z jego większym przystosowaniem, pozwala to w rezultacie korygować funkcję celu, by na przykład bardziej preferowała mniejsze sieci,
  \item maksymalna liczba zrównolegleń - liczba całkowita dodatnia,
  \item liczba obliczań przystosowania - liczba całkowita dodatnia,
  \item flaga wyświetlania postępu - wartość prawda/fałsz,
  \item flaga rysowania wykresu - wartość prawda/fałsz.
\end{nospitemize}
Dane treningowe, walidacyjne i testowe mogą mieć wyjście klasyfikacyjne lub regresyjne, a wejście regresyjne, jednakże muszą być to jednolite względem siebie oraz muszą mieć tą samą ilość kolumn z atrybutami wejściowymi jak i wyjściowymi.
Funkcja ma zakończyć działanie po wykonaniu zadanej liczby generacji. Funkcja ma udostępniać pod wybraną przez uczestnika lokalizacją: dane wszystkich osobników w sposób umożliwiający odtworzenie ich, ich nauczone wagi oraz wartości przystosowania przed przejściem przez funkcję korygującą.
Po wywołaniu funkcji przeprowadzony jest pełny przebieg ewolucji pokazany na \cref{fig:genetic_flow}. Jest tworzona populacja początkowa, po czym dla każdego osobnika jest $n$-krotnie ($n$=liczba obliczań przystosowania) obliczana funkcja przystosowania, w ramach której jest tworzona i uczona sieć neuronowa, a wyniki funkcji są uśredniane by uzyskać końcowe przystosowanie. Następnie wykonywane są kolejno selekcja, krzyżowanie, mutacja i wyliczenie przystosowania nowej populacji (ponieważ wszystkie nowe osobniki mają być wzięte do populacji w nowej generacji), kroki podane w niniejszym zdaniu są powtarzane aż do osiągnięcia warunku końcowego będącego dojściem do zadanej (lub domyślnej) liczby generacji.

Druga funkcja ma przyjomwać osobnika, liczbę atrybutów wejściowych i wyjściowych oraz flagę czy dane wyjściowe są klasyfikacyjne. Ich format to:
\begin{nospitemize}
  \item osobnik - osobnik deterministycznie określający strukturę sieci neuronowej,
  \item liczba atrybutów wejściowych - liczba całkowita dodatnia,
  \item liczba atrybutów wyjściowych - liczba całkowita dodatnia.
  \item flaga klasyfikacji - wartość prawda/fałsz
\end{nospitemize}
Funkcja zwraca stworzony model sieci oraz parametry z jakimi ma być uczona: wielkość wsadu i liczba epok. Mają mieć one format:
\begin{nospitemize}
  \item model - sieć umożliwiająca uczenie, zapisanie i odtworzenie wag oraz danie odpowiedzi do danych,
  \item wielkość wsadu - liczba całkowita dodatnia
  \item liczba epok - liczba całkowita dodatnia
\end{nospitemize}
W ramach swojego działania funkcja odczytuje parametry osobnika i tworzy na ich podstawie sieć.